\section{Historical three-family main comparisons}
\label{app:historical-main}
These computational results are reproduced from the source manuscript and retain its original
population, checkpoints, methods and denominators. They are not 22-family results or a matched
old-versus-new comparison. The new bridge remains pending in Table~\ref{tab:22-controls}.

\subsection{Evaluation protocol and matched budgets}
We compared correctly conditioned, shared-null and cyclically mislabelled models using the same
architecture, optimization budget, attempt count and final-checkpoint rule across three independent
training seeds. Component-disjoint held-out products were excluded from training and checkpoint
selection. Verified exact-L1 yield counts attempts passing molecular validation, decomposition and
exact forward replay, with \emph{all} attempts in the denominator. Post-hoc rejection cannot increase
the proposal's per-attempt acceptance probability (Proposition~\ref{prop:posthoc-budget}). We report
means and sample standard deviations across training seeds, distinguishing retraining variation from
conditional Monte Carlo variation. Appendix~\ref{app:evaluation-protocol} specifies the comparisons,
coverage and precision metrics, and matched budgets.

\subsection{Shared generation achieved exact assembly across three programs}
We first tested whether one conditional flow could learn all three reactions while preserving Ugi
performance. Across
\ForgeProseAttemptsPerProgram{} held-out attempts per program and seed, FORGE achieved verified
exact-L1 yields of \ForgeTableOneConditionedUgiYield\% for Ugi,
\ForgeTableOneConditionedBLYield\% for repeated aza-Michael addition and
\ForgeTableOneConditionedLXYield\% for repeated reductive amination
(Table~\ref{tab:iclr-shared-programs}). Every verifier-returned trace passed exact forward replay in
every program and seed, as required by construction. Reductive amination showed greater between-seed
variation and higher decomposition ambiguity, indicating less consistent learning in this lower-data
setting.

FORGE produced no fixed-coordinate violations or support overflows across
\ForgeProseConditionedTotalAttempts{} held-out attempts. These runs established exact assembly across
all three programs, with the highest yield for Ugi.

\begin{table}[H]
\caption{\textbf{Exact assembly under three ionizable-lipid reaction programs.}
Exact-L1 yields are mean $\pm$ sample standard deviation across three seeds. The final column reports
decomposition coverage and the forward-replay verifier invariant. Per-seed values:
Appendix~\ref{app:results}.}
\label{tab:iclr-shared-programs}
\centering
\small
\setlength{\tabcolsep}{4.6pt}
\begin{tabular}{@{}lcccc@{}}
\toprule
Program & FORGE conditioned & Shared null & Cyclic program &
\shortstack{Decomp. coverage /\\ replay invariant (\%)} \\
\midrule
Ugi & $96.4\pm2.1$ & $34.8\pm7.1$ & $93.8\pm4.2$ & 99.9/100.0 \\
Aza-Michael & $72.2\pm1.6$ & $19.2\pm4.1$ & $71.5\pm2.2$ & 100.0/100.0 \\
Reductive amination & $52.3\pm3.4$ & $37.0\pm2.2$ & $50.7\pm4.9$ & 86.6/100.0 \\
%
\end{tabular}
\end{table}

\subsection{Program conditioning improved exact-assembly yield}
To isolate the effect of reaction information during generation, we compared the conditioned
model with the shared-null and cyclic arms under the same attempt budget and post-hoc verifier.
Shared generation without program information yielded \ForgeTableOneNullUgiYield\%,
\ForgeTableOneNullBLYield\% and \ForgeTableOneNullLXYield\% exact L1 for Ugi, aza-Michael addition
and reductive amination, respectively (Table~\ref{tab:iclr-shared-programs}). For Ugi, conditioning
increased yield by \ForgeTableOneNullUgiDifferencePP{} percentage points (three-seed range
\ForgeTableOneNullUgiRangeLowPP{} to \ForgeTableOneNullUgiRangeHighPP{} points).

Assigning the wrong reaction identity reduced exact-L1 yield to
\ForgeTableOneCyclicUgiYield\%, \ForgeTableOneCyclicBLYield\% and
\ForgeTableOneCyclicLXYield\% across the three programs. The cyclic control retained precursor roles, reaction-core
coordinates and program depth, isolating the contribution of the reaction-identity label within a
structured program. The larger conditioned-versus-null differences supported a benefit from program
information during generation. The small conditioned-versus-cyclic differences did not establish
a strong independent contribution from the identity label.

\subsection{Decoder and architecture effects differed across programs}
With seed-0 weights fixed, enforcing Ugi core hydrogen states increased exact-L1 yield with nearly
unchanged molecular validity; changing source marginals from program-role to global reduced yield.
The core-saturation constraint changed the generated graph, with the same verifier applied to both
decoders (Appendix~\ref{app:ablation-results}).

A separate three-seed study compared joint denoising with parameter-matched independent precursor
roles (FACT). Its full Transformer was trained independently of the production model. Mean exact-L1
products per 1,000 attempts were \ForgeProseMechanismFullUgi{} versus
\ForgeProseMechanismFactMatchedUgi{} for Ugi and \ForgeProseMechanismFullBL{} versus
\ForgeProseMechanismFactMatchedBL{} for aza-Michael addition. Reductive amination favored FACT:
\ForgeProseMechanismFullLX{} versus \ForgeProseMechanismFactMatchedLX{}. Removing cross-attention reduced Ugi and reductive-amination yield but
increased aza-Michael yield; removing gradient-conflict control slightly improved all three means.
These ablations support reaction-dependent effects, not a uniform benefit from every module.

\subsection{Generated precursors extended beyond the training catalogues}
We next asked whether FORGE could preserve exact assembly while generating precursor constitutions
absent from training. Under the common Ugi benchmark, the finite-inventory methods produced
\ForgeProseCommonSelectorExactPerThousand{} verified exact-L1 products per 1,000 attempts but zero
component-novel products by construction. FORGE produced
\ForgeProseCommonForgeExactPerThousand{} verified products and
\ForgeProseCommonForgeDistinctPerThousand{} distinct component-novel products per 1,000 attempts
(Table~\ref{tab:ugi-common-benchmark}).

\begin{table}[htbp]
\caption{\textbf{Exact Ugi assembly with precursor constitutions absent from
method-visible training.} Values are mean $\pm$ sample standard deviation per 1,000 attempts across
three seeds under a common component-disjoint protocol and verifier. Component-novel is
verifier-relative and does not establish strong product-level catalogue escape. Whole-molecule
baselines are assessed post hoc. $\dagger$, structural zero; N/E, not estimable.}
\label{tab:ugi-common-benchmark}
\centering
\scriptsize
\setlength{\tabcolsep}{2.2pt}
% Column 1 is fixed so the numeric columns, not the model names, decide the table's natural
% width. The generated rows separate the three baseline groups with a rule and no heading.
% 3.3cm is the widest model name ("Learned inventory selector", 92.4pt) plus a little; the
% \resizebox is a guard, since the header arrows push the natural width just past the block.
\resizebox{\textwidth}{!}{%
\begin{tabular}{@{}>{\raggedright\arraybackslash}p{3.3cm}cccccc@{}}
\toprule
Models & \shortstack{Valid / 1,000\\$\uparrow$} &
\shortstack{Verified L1 / 1,000\\$\uparrow$} &
\shortstack{Distinct L1 / 1,000\\$\uparrow$} &
\shortstack{Comp.-novel / 1,000\\$\uparrow$} &
\shortstack{Held-comp. / 1,000\\$\uparrow$} & Diversity \\
\midrule
RGFN & $0.0\pm0.0$ & $0.0\pm0.0$ & $0.0\pm0.0$ & $0.0\pm0.0^{\dagger}$ & $0.0\pm0.0$ & N/E \\
\midrule
DeFoG unconditional & $346.5\pm42.3$ & $280.6\pm26.6$ & $280.6\pm26.6$ & $280.6\pm26.6$ & $15.5\pm2.4$ & $0.727\pm0.014$ \\
GenMol/SAFE & $627.5\pm6.2$ & $0.0\pm0.0$ & $0.0\pm0.0$ & $0.0\pm0.0$ & $0.0\pm0.0$ & $0.873\pm0.009$ \\
\midrule
Learned inventory selector & $1000.0\pm0.0$ & $1000.0\pm0.0$ & $849.5\pm13.7$ & $0.0\pm0.0^{\dagger}$ & $0.0\pm0.0$ & $0.717\pm0.001$ \\
Finite catalogue oracle & $1000.0\pm0.0$ & $1000.0\pm0.0$ & $837.6\pm8.2$ & $0.0\pm0.0^{\dagger}$ & $0.0\pm0.0$ & $0.719\pm0.000$ \\
Shared-null + post-hoc & $194.7\pm15.6$ & $114.4\pm10.9$ & $108.5\pm9.7$ & $108.5\pm9.7$ & $0.4\pm0.2$ & $0.754\pm0.014$ \\
FACT-matched & $60.8\pm20.2$ & $52.0\pm15.6$ & $50.8\pm14.4$ & $50.8\pm14.4$ & $0.0\pm0.0$ & $0.667\pm0.009$ \\
FACT-generous & $500.7\pm213.8$ & $296.7\pm186.0$ & $288.3\pm175.3$ & $288.3\pm175.3$ & $0.1\pm0.2$ & $0.774\pm0.007$ \\
\midrule
\cellcolor{forgerow}FORGE &\cellcolor{forgerow} $965.0\pm21.9$ &\cellcolor{forgerow} $963.9\pm21.4$ &\cellcolor{forgerow} $862.6\pm19.2$ &\cellcolor{forgerow} $862.6\pm19.2$ &\cellcolor{forgerow} $1.4\pm1.7$ &\cellcolor{forgerow} $0.643\pm0.022$ \\
%
\end{tabular}
}
\end{table}

Against the fixed training catalogues, component-novel verified exact-L1 yields reached
\ForgeProseCatalogueForgeUgiComponentNovelPerThousand{},
\ForgeProseCatalogueForgeBLComponentNovelPerThousand{} and
\ForgeProseCatalogueForgeLXComponentNovelPerThousand{} products per 1,000 attempts across Ugi,
aza-Michael addition and reductive amination. The corresponding finite catalogues had zero
component-novel yield by construction (Appendix~\ref{app:novelty-results}). Representative Ugi
products and exact precursors appear in Appendix~\ref{app:generated-examples}.

Among the external baselines, DeFoG produced \ForgeProseCommonDefogExactPerThousand{} verified
products and recovered held components more often than FORGE
(\ForgeProseCommonDefogHeldPerThousand{} versus \ForgeProseCommonForgeHeldPerThousand{} per 1,000
attempts). GenMol/SAFE produced no exact Ugi-L1 product, and RGFN emitted no admissible output. FORGE therefore combined reaction consistency with verifier-relative component novelty,
although DeFoG recovered the reserved components more often. Product-level catalogue escape
remained unmeasured.

We also compared the generated molecules with known lipids excluded from training. For every
method, the two sets differed in their overall structural features.
Precision and coverage in fingerprint and descriptor spaces, together with a source-study-grouped classifier
two-sample test, found that no method reproduced the full observed-lipid distribution
(Appendix~\ref{app:novelty-results}). The measured expansion concerns reaction-consistent identities;
it does not establish biological utility.

